% Equation bank for OFDM 16-QAM project
% Copy into the IEEE report as needed.

\subsection{Constelação 16-QAM}
Para a modulação 16-QAM, cada símbolo transporta
\begin{equation}
    b = \log_2(16)=4
\end{equation}
bits. A constelação quadrada não normalizada utiliza níveis $I,Q\in\{-3,-1,+1,+3\}$. A energia média é $E_s=10$, de modo que a constelação normalizada é dada por
\begin{equation}
    s_m = \frac{I_m+jQ_m}{\sqrt{10}}, \qquad \mathbb{E}\{|s_m|^2\}=1.
\end{equation}

\subsection{Canal discreto}
O canal multipercurso especificado é
\begin{equation}
    h[n]=[0.3,\,-0.5,\,0,\,1,\,0.2,\,-0.3]^T.
\end{equation}
Sua resposta em frequência para $N=32$ subportadoras é
\begin{equation}
    H[k]=\sum_{n=0}^{5} h[n]e^{-j2\pi kn/32}, \qquad k=0,1,\ldots,31.
\end{equation}

\subsection{Prefixo cíclico}
Como o canal possui comprimento $L_h=6$, sua memória é $L_h-1=5$. Assim, o prefixo cíclico mínimo é
\begin{equation}
    N_{CP,min}=L_h-1=5.
\end{equation}

\subsection{Modelo OFDM no domínio da frequência}
Com prefixo cíclico suficiente, a convolução linear do canal equivale a uma convolução circular no bloco útil. Após a DFT, cada subportadora obedece a
\begin{equation}
    Y[k]=H[k]X[k]+V[k].
\end{equation}

\subsection{Equalização ZF}
O equalizador zero-forcing é
\begin{equation}
    E[k]=\frac{1}{H[k]},
\end{equation}
e a saída equalizada é
\begin{equation}
    \hat{X}[k]=\frac{Y[k]}{H[k]}=X[k]+\frac{V[k]}{H[k]}.
\end{equation}
A variância do ruído após equalização é, portanto,
\begin{equation}
    \sigma_{eq,k}^{2}=\frac{\sigma_V^2}{|H[k]|^2}.
\end{equation}

\subsection{Normalização de ruído}
Para a convenção de SNR média no domínio da frequência antes da equalização,
\begin{equation}
    \sigma_V^2=\frac{\frac{1}{N}\sum_{k=0}^{N-1}|H[k]|^2}{\mathrm{SNR}_{lin}}=\frac{1.47}{\mathrm{SNR}_{lin}}.
\end{equation}
Com a convenção da FFT do NumPy, se o ruído for somado no domínio do tempo,
\begin{equation}
    \sigma_t^2=\frac{\sigma_V^2}{N}=\frac{1.47}{32\,\mathrm{SNR}_{lin}}.
\end{equation}

\subsection{SER}
A taxa de erro de símbolo da subportadora $k$ é estimada por
\begin{equation}
    \mathrm{SER}_k=\frac{N_{err,k}}{N_{sym,k}}.
\end{equation}
A SER média global é
\begin{equation}
    \mathrm{SER}_{avg}=\frac{\sum_k N_{err,k}}{\sum_k N_{sym,k}}.
\end{equation}
